\begin{table}[htbp]
  \centering
  \caption{Parameters of the simulation model. One time step
  equals one hour.}
  \label{tab:simulator_event_parameters}
  \scriptsize
  \renewcommand{\arraystretch}{0.92}
  \setlength{\tabcolsep}{3pt}
  \begin{tabular}{@{}p{0.35\linewidth}p{0.20\linewidth}p{0.38\linewidth}@{}}
    \toprule
    Component & Value & Role \\
    \midrule
    Reference distributions from Antarctic AWS stations
      & Panda100, Panda200, Taishan
      & Phase-randomized Antarctic AWS series with empirical marginal matching. \\
    Sequence length and interval
      & 70,000; 1 h
      & One independently generated simulated ground-truth time series per seed. \\
    Simulated blowing-snow event process
      & alternating semi-Markov storm/calm phases
      & Lognormal phase duration with log-scale standard deviation 0.35 and a
        minimum duration of 24 time steps. \\
    Storm / within-storm event duty
      & 0.65 / 0.78
      & Derived from the event-coverage target after clipping. \\
    Median storm / calm duration
      & 1,440 / 775 time steps
      & Parameters of the alternating calm/storm process. \\
    Event pulse duration / gap
      & Uniform 28--80 / 8--10 time steps
      & Bursty blowing-snow episodes within storm phases. \\
    Wind thresholds and event margin
      & 8 / 12 / 1.4 m s$^{-1}$
      & Weak threshold, strong threshold, and event wind-floor margin. \\
    Hysteresis on / off
      & 0.60 / 0.30
      & Converts candidate pulses and wind credibility into event intervals. \\
    Particle / flux / thermal probability
      & 0.40 / 0.20 / 0.40
      & One event type is sampled for each complete event run. \\
    Event-type process smoothing
      & 0.65
      & Low-pass coefficient for event-type-specific simulated processes. \\
    Target-effect lag
      & 6 time steps
      & Delay from the event-type process to target effects. \\
    Particle process scales
      & 0.62 mm; 8.2 m s$^{-1}$
      & Diameter and particle-velocity effects. \\
    Flux process map
      & scale 0.0014; offset 1.25; clip 2.4
      & Linear event-type effect before clipping. \\
    Thermal process scale
      & 6.4 $^{\circ}$C
      & Snow-surface-temperature effect. \\
    Warning-signal lead / noise
      & 16 time steps / 0.01
      & Linear pre-event ramp plus Gaussian noise, clipped to $[0,1]$. \\
    \bottomrule
  \end{tabular}
\end{table}
